% ═══════════════════════════════════════════════════════════════════════════════
% APÉNDICE B: MATRIZ COMPARATIVA - INSTITUTO ATENEO DE COLIMA
% Estudio Documental: Corrupción en México (2006-2012)
% ═══════════════════════════════════════════════════════════════════════════════

\nota{
  Este apéndice presenta la matriz comparativa utilizada para integrar 
  evidencia cualitativa (mecanismos institucionales) con indicadores 
  cuantitativos anidados (montos, víctimas, sentencias). La matriz 
  permite generar metainferencias trazables por caso.
}

\section{Matriz Comparativa de Variables, Dimensiones e Indicadores por Caso}

\begin{landscape}
\begin{table}[H]
\centering
\caption{Matriz Comparativa por Caso (Apéndice B)}
\label{tab:matriz_comparativa}
\tiny
\begin{tabular}{@{}p{2cm}p{2.5cm}p{3.5cm}p{4cm}p{6.5cm}@{}}
\toprule
\textbf{Caso y tipo} & \textbf{Dimensión analítica / factor asociado} & \textbf{Mecanismo institucional (CUAL)} & \textbf{Indicadores cuantitativos (cuan)} & \textbf{Fuentes trianguladas (APA 7)} \\
\midrule
Caso 1 (Principal) \newline Genaro García Luna & Concentración de mando, presupuesto e inteligencia en SSP sin contrapesos. & Captura estatal y macrocriminalidad: cohecho con Cártel de Sinaloa + red offshore Weinberg. & PF: 5,000 a >35,000; 30 contratos (>727 MDD y 528 mdp); prisión >38 años (2024); civiles \textasciitilde2,500 MDD (2025) y >578 MDD (2026). & DOJ (2024, exp. 19-CR-576); UIF (2024, 2025, 2026); Astorga (2015); Vázquez Valencia (2019); Infobae (2024). \\
\addlinespace
Caso 2 (Principal) \newline Refinería Bicentenario & Anuncio político prematuro, preinversión débil y sin candados transexenales. & Opacidad y mala planeación en megainfraestructura con transferencia de pasivos. & Proyección >9,000 MDD; gasto Pemex >9,000 mdp; barda >500 mdp; deuda Hidalgo \textasciitilde1,500 mdp liquidada sep. 2025; 195.5 mdp en 2020. & ASF (2015, 2016); Cancino (2020); Sandoval Ballesteros (2024); El Financiero (2025). \\
\addlinespace
Caso 3 (Principal) \newline Guardería ABC & Desregulación de estancias IMSS, conflicto de interés y omisión de inspecciones. & Subrogación desregulada con dilución de la responsabilidad administrativa estatal. & Subrogadas: 757 a 1,374 (+81.5\,\%) vs. 142 directas; 14 de 40 empleados; 49 fallecidos y >100 lesionados graves. & CNDH (2009); SCJN (2009, 2010, 2025); Núñez Esquer (2026); CIDH (2020); Serna de la Garza (2010); Osorno (2010). \\
\addlinespace
Contraste 1 \newline Sector Salud Veracruz & Saqueo presupuestal y empresas fantasma; encuadre viral impreciso. & Contraste entre rumor mediático («agua destilada») y desvío hospitalario acreditado. & 0 evidencia agua destilada (COFEPRIS); 46,984 pruebas VIH y 16.8 t fármacos caducos; 2,933 mdp obra; >3,617 mdp fantasma. & Secretaría de Salud / COFEPRIS (2017); Connectas (s.\,f.); Animal Político (s.\,f.); Santiago y Morales (2025); Ángel (2025); Sola-Morales y Zurbano-Berenguer (2019). \\
\addlinespace
Contraste 2 \newline Refinería Olmeca / Dos Bocas & Sobrecostos en megainfraestructura y polarización en eslóganes partidistas. & Contraste entre consigna viral («no refina») y desempeño técnico-financiero real. & Costo: 8,000 a >21,000 MDD (\textasciitilde2.6x); producción 212--263 mil bd (fin 2025) con caída a 141 mil bd (jun. 2026). & Alegría (2026); Martínez Escamilla (2026); El Universal / Celaya (2026); Cancino (2020); Reyes Heroles (2003). \\
\bottomrule
\end{tabular}
\end{table}
\end{landscape}